\chapter{Conclusions and Future Work}
\label{ch:conclusoes}

\section{Main Conclusions}

\rat\ meets the objectives defined in Chapter~\ref{ch:intro}: it delivers integrated static and
dynamic RAT analysis with explainable outputs rather than opaque verdicts. Evaluation
(Section~\ref{sec:discussion-results}) shows internally consistent scoring on synthetic profiles and
zero false positives on a small benign corpus, while explicitly not claiming field-perfect
detection.

The \textbf{static pipeline} implements seven analysis phases backed by extensive automated backend
tests. The \textbf{Hyper-V sandbox} provides two orchestration paths that share snapshot-isolated
execution. \textbf{Interfaces} include a web application with streaming progress, tri-panel
reporting, and an optional Gemini assistant, together with a WPF desktop entry point. The project
includes \textbf{220 automated tests} across Python, TypeScript, and .NET stacks, complemented by
\textbf{architecture diagrams}, a security guide, and deployment manuals. Throughout, the design
emphasises \textbf{explainability}: a factorised risk score, decompiled views, and optional
LLM-assisted interpretation of pseudo-C excerpts.

\section{Limitations}

Several limitations bound the claims that can be made from the current evaluation. Validation
relies on eight benign PEs, 59 synthetic profiles, and public IOC hashes---not on real RAT binaries
executed on the development host (see Section~\ref{sec:discussion-results}). Perfect synthetic
accuracy verifies specification consistency, not field detection rate; construct and spectrum bias
therefore apply. The MalwareBazaar batch script was implemented but not executed under the
host-safety policy adopted for this project. Dynamic analysis was exercised only with a benign
smoke-test binary, and Path~A telemetry remains shallower than Path~B. On the web/API,
\texttt{.cs} uploads receive string heuristics only. Heuristic scoring may false-positive on
legitimate network-administration tools, and Hyper-V ties deployment to Windows-capable laboratory
hosts.

\section{Future Work}

Natural extensions follow from the modular architecture. \textbf{Multi-VM and queue scaling} would
deploy parallel sandboxes with worker pools and snapshot pools to raise throughput beyond
single-run serialisation. \textbf{Linux malware support} could extend dynamic analysis to ELF
specimens via KVM guests alongside the existing Windows Hyper-V path. \textbf{Automated behavioural
clustering} would group sandbox event traces into similarity classes, surfacing campaign-level
patterns without manual log review. \textbf{Memory forensics}---for example Volatility-style
inspection of guest RAM snapshots---could recover injected code and unpacked payloads missed by
file-system telemetry alone. \textbf{MITRE ATT\&CK mapping} would align static and dynamic
indicators with standard tactic/technique identifiers for reporting and curriculum alignment.

Further improvements include IDA Pro integration for deeper native analysis beyond Ghidra
pseudo-C; automatic C2 configuration extraction from obfuscated samples; enriched dynamic telemetry
(Sysmon, ETW) in the HTTP vm-agent path to match Path~B depth; support for macros, scripts, and
mobile (APK) formats where institutional policy allows; a curated malware hash database
(MalwareBazaar/MALPEDIA) with supervised static batch runs; C2 network traffic capture on an
instrumented internal VLAN; a backend LLM proxy with additional models and automatic
suspicious-function explanations; and YARA rule management in the web interface for classroom rule
authoring.

\rat\ provides a solid, extensible foundation for academic malware analysis and future exploration.
